\section{A classical ordinary CM input with the required normalization}
\label{ct-section}

The integrality theorem for reduced CM theta functions gives a shorter
route to part of the ordinary theta input.  We state precisely the
classical result used, match the even theta function and the period,
and prove the stronger norm assertion for finite combinations.  The
latter does not follow merely by bounding each theta function
separately.

Fix throughout an abstract field embedding $\iota:\mathbb C\hookrightarrow
\mathbb C_p$, compatible with the chosen embedding of the algebraic
coefficients.  The submitted construction allows an abstract field
isomorphism, which is more than is needed here.  Every transported
Taylor series, torsion character and complex period below uses this
same $\iota$.  No continuity for the complex topology, and no natural
embedding between the two fields, is asserted.

\subsection{The classical input and its scope}

We use the following classical facts in their elliptic-curve form.
For a good ordinary CM elliptic curve at a split prime $p\ge5$,
the reduced theta function with divisor $[O]$ and derivative one
has integral formal expansions after prime-to-$p$ Mumford translation;
the constant term is a unit unless the translated divisor contains
$O$.  The formal group has an integral isomorphism with
$\widehat{\mathbb G}_m$, whose period relative to a good invariant
differential is a unit.  The precise references are
Bannai--Kobayashi, Proposition~2.15, Corollary~2.16(i), and
\S3.1, in the version specified in the bibliography
\cite{bannai-kobayashi-2007}.

These statements use a good integral model, its invariant
differential, and reduced theta functions; arbitrary exponential
renormalizations are not included.  Prime-to-$p$ torsion includes
points of even order because $p\ge5$.  The full CM ring is the
maximal order.  The associated Hermitian form must be compatible
with complex multiplication; in one dimension it is a real
multiple of $z\bar w$, so this condition is automatic.
We import these classical theorems explicitly.  Their proofs are
not represented here as consequences of the finite moment argument.

Let $d$ be a positive squarefree integer with $d\equiv7\pmod8$,
put
\[
 k=\mathbb Q(\sqrt{-d}),\qquad
 \tau=\frac{1+i\sqrt d}{2},\qquad
 \Lambda=\mathcal O_k=\mathbb Z+\mathbb Z\tau,
 \qquad H(z,w)=\frac{2}{\sqrt d}z\bar w,
 \tag{CT1}
\]
and write $E(z,w)=\operatorname{Im}H(z,w)$.
For a function $F$ and $t\in\mathbb C$, define
\[
 (U_tF)(z)=
  e^{-\pi H(z,t)-\pi H(t,t)/2}F(z+t).
 \tag{CT2}
\]
Thus the Hermitian form is linear in its first variable.
Direct multiplication of the exponentials gives
\[
                   U_tU_s=e^{-\pi iE(t,s)}U_{t+s}.
 \tag{CT3}
\]
For rational lattice points $s,t\in\Lambda\otimes\mathbb Q$,
the factor in (CT3) is a root of unity, because
$E(\Lambda,\Lambda)\subset\mathbb Z$.

Fix a good integral CM model $\mathscr E$ over the valuation ring
$R$ of a finite extension of $\mathbb Q_p$, and a complex
uniformization compatible with a good invariant differential.
After a finite extension, take the full CM action to be defined
on this elliptic scheme.
If its period lattice is $\Gamma=c\Lambda$, its Hermitian form
satisfies
\[
                        H_\Gamma(cz,cw)=H(z,w).
 \tag{CT4}
\]
All the functions and torsion points used below can be defined
after a further finite extension.  We then embed the coefficients
in $\mathbb C_p$.  The chosen good model and its multiplicative
formal-group coordinate are part of the classical input.
The multiplicative coordinate may be defined over the completed
maximal unramified extension, or over $\mathcal O_{\mathbb C_p}$;
it need not be defined over the finite extension with valuation
ring $R$.  The finite ring $R$ will be used for the curve, its
line bundles and their finite free modules of global sections.
Integral changes of formal coordinate are then made after scalar
extension to $\mathbb C_p$.

\subsection{Identification of the even theta function}

Let $\vartheta_\Gamma$ be the reduced odd theta function for
$\Gamma$, with simple zeros exactly at $\Gamma$ and
$\vartheta'_\Gamma(0)=1$.  Put
$\vartheta(z)=\vartheta_\Gamma(cz)$ and $h=1/2$, and define
\[
                 \theta(z)=\frac{(U_h\vartheta)(z)}
                                      {(U_h\vartheta)(0)}.
 \tag{CT5}
\]
The denominator is nonzero because $h\notin\Lambda$.

\begin{lemma}\label{ct-even-identification}
The function (CT5) is entire and even, satisfies $\theta(0)=1$,
has simple zeros exactly at $h+\Lambda$, and satisfies
\[
 \theta(z+\ell)=(-1)^{N_{k/\mathbb Q}(\ell)}
            e^{\pi H(z,\ell)+\pi H(\ell,\ell)/2}\theta(z)
                   \qquad(\ell\in\Lambda).
 \tag{CT6}
\]
It is the unique entire function with this transformation law,
constant term one, and no zeros outside $h+\Lambda$.
\end{lemma}
\begin{proof}
The odd reduced theta function has semicharacter
$\alpha(u+v\tau)=(-1)^{u+v+uv}$.
Translation by $h$ multiplies it by
$e^{2\pi iE(h,u+v\tau)}=(-1)^v$, since
$E(h,u+v\tau)=-v/2$.  The translated semicharacter is therefore
$(-1)^{u+uv}$.  On the other hand
\[
 N(u+v\tau)=u^2+uv+\frac{d+1}{4}v^2\equiv u+uv\pmod2,
\]
because $(d+1)/4$ is even.  This proves (CT6).
Translation gives the stated zeros, and division by the value at
zero normalizes the constant term.

For completeness, every nonzero entire function with this
multiplier has exactly one zero, with multiplicity, in a
fundamental parallelogram.  Choose its boundary to avoid zeros.
The logarithmic derivative changes on translation by $\ell$ by
$\pi H(1,\ell)$.  Pairing the two horizontal and the two
vertical edges gives the boundary integral
\[
       \frac{\pi}{\operatorname{Im}\tau}(\tau-\bar\tau)
                              =2\pi i.
\]
The argument principle gives the claim.  Any other function in
the stated class consequently has one simple zero in the same
coset.  Its quotient by (CT5) is entire and $\Lambda$-periodic,
hence bounded and constant, and the value at zero fixes that
constant.  Finally $\theta(-z)$ has the same multiplier, the same
zero coset, and the same value at zero.  Uniqueness proves
evenness.
\end{proof}

Since $N(\ell)=|\ell|^2$ for $\ell\in k$, (CT6) is equivalent
to
\[
               e^{\pi i|\ell|^2}U_\ell\theta=\theta.
 \tag{CT7}
\]
Thus this is precisely the even theta normalization used in the
ordinary moment construction.  The condition
$d\equiv7\pmod{16}$ in that construction is stronger than the
congruence needed in this identification.

\subsection{One common unit period}

Let $\lambda(u)$ be the formal logarithm for the good invariant
differential on $\mathscr E$, normalized by $\lambda'(0)=1$ in
an integral parameter $u$ whose linear differential is that basis.
Choose the integral multiplicative isomorphism in the form
\[
                T=\exp(\lambda(u)/\Omega_p)-1,
                  \qquad |\Omega_p|_p=1.
 \tag{CT8}
\]
Its inverse is integral with unit linear coefficient.  In the
original variable $z$ for $\Lambda$, the same coordinate reads
\[
              z=\omega\log(1+T),\qquad
                        \omega=\Omega_p/\iota(c).
 \tag{CT9}
\]
This identity is an identity of transported formal Taylor series;
it is not a topological identification of complex and $p$-adic
neighborhoods.

Let $g_2,g_3$ be the Weierstrass invariants of $\Lambda$ and choose
$\delta\ne0$ with
\[
                       \delta^{12}=g_2^3-27g_3^2.
 \tag{CT10}
\]
The good model can be general Weierstrass.  With its invariants
$c_4,c_6,\Delta$, the change of differential gives
$c_4/12=\iota(g_2/c^4)$ and
$c_6/216=\iota(g_3/c^6)$.  Hence the identity
$c_4^3-c_6^2=1728\Delta$ gives
\[
               \iota(\delta)^{12}=\iota(c)^{12}\Delta.
 \tag{CT11}
\]
Since $|\Delta|_p=1$, one has
$|\iota(\delta)|_p=|\iota(c)|_p$.  Therefore
\[
             \Omega=\iota(\delta)\omega,
                    \qquad |\Omega|_p=1,
                    \qquad\omega=\Omega/\iota(\delta).
 \tag{CT12}
\]
The same $\Omega$ works for every torsion translate and every
finite family of directions: no period is chosen after selecting
the translate or a derivative order.

\begin{proposition}\label{ct-translates-integral}
For every $t\in\Lambda\otimes\mathbb Q$ of order prime to $p$
modulo $\Lambda$, the series
\[
                 F_t(T)=\widehat{U_t\theta}
                                  (\omega\log(1+T))
 \tag{CT13}
\]
is integral.  If $t+h\notin\Lambda$, it is a unit series.
\end{proposition}
\begin{proof}
For the good model, the classical input applies to the point $ch$
of order two, so the formal expansion of $U_h\vartheta$ has unit
constant term.  The division in (CT5) is therefore by a $p$-adic
unit.  By (CT3), $U_t\theta$ is this same unit inverse times a
root of unity times $U_{t+h}\vartheta$.
The point $c(t+h)$ still has order prime to $p$.
The classical theorem makes its reduced-theta formal expansion
integral, and a unit unless $t+h\in\Lambda$.  Substitution of
the integral inverse of (CT8) preserves integrality and unit
constant terms.  The change from $cz$ to $z$ is exactly (CT9),
which proves the displayed assertion.
\end{proof}

Proposition~\ref{ot-logarithm} consequently supplies norm-at-most-one
measures for all these one-dimensional period-scaled Taylor
functionals.  The exceptional translate has a zero constant term,
but its integrality and its measure remain valid.

\subsection{An integral expansion detects all global sections}

\begin{lemma}\label{ct-expansion-saturated}
Let $\mathscr X/R$ be a smooth proper curve with geometrically
integral special fiber, let $x\in\mathscr X(R)$, and let
$\mathscr L$ be a line bundle.  Choose an integral trivialization
of $\mathscr L$ in the formal neighborhood of $x$ and an integral
parameter there.  Then the expansion map
\[
 H^0(\mathscr X,\mathscr L)\longrightarrow R[[u]]
 \tag{CT14}
\]
is injective and is saturated: a global section over the fraction
field has integral expansion if and only if it is integral.
After scalar extension to $\mathbb C_p$, the coefficient
supremum norm of an expansion equals the maximum absolute value
of its coordinates in any $R$-basis of $H^0(\mathscr X,\mathscr L)$.
The same assertion holds if the expansion is multiplied by an
integral unit series, or its variable is changed by an integral
series with unit linear coefficient.
\end{lemma}
\begin{proof}
The module of global sections is finite over $R$ by properness
and torsion-free because the line bundle is flat.  It is therefore
free.  Choose an $R$-basis $s_1,\ldots,s_q$.
Reduction modulo a uniformizer injects this module modulo that
uniformizer into the space of sections on the special fiber:
this is the initial part of the cohomology sequence of
$0\to\mathscr L\xrightarrow{\pi}\mathscr L
\to\mathscr L|_{\mathscr X_k}\to0$.
On an integral smooth curve, a section which vanishes in the
completed local ring at a point vanishes globally.  Indeed its
germ vanishes in the local discrete valuation ring, hence on a
nonempty open set, and the generic value then vanishes.  The
local ring embeds in its power-series completion, since a
nonzero element has finite valuation.  Thus the reductions of
the expansions of the $s_i$ are linearly independent.

Some finite list of coefficient functionals already detects this
independence: their kernels form a decreasing sequence of
subspaces in a finite-dimensional vector space, whose full
intersection is zero.  Consequently the coefficient matrix has
a $q\times q$ minor which is a unit of $R$.
For any coefficients $a_i\in\mathbb C_p$, integrality of every
expansion coefficient of $\sum a_i s_i$ implies integrality of
the selected finite list and hence of all $a_i$, by the integral
inverse of that minor.  More generally the same inverse has
entries of absolute value at most one, so
\[
 \max_i|a_i|_p\le
   \sup_n|[u^n]\textstyle\sum_i a_i s_i|_p
                  \le\max_i|a_i|_p.
\]
This proves all the assertions about saturation and norms,
including injectivity.  Multiplication by an integral unit has
integral inverse, and so does a change of variable with unit
linear coefficient.  Applying each operation and its inverse
gives equality of norms.  The zero-dimensional space of sections
causes no exception.
\end{proof}

\begin{lemma}\label{ct-torsion-isometry}
In the situation of Lemma~\ref{ct-expansion-saturated}, let
$\mathscr X$ be an elliptic scheme.  Let $v$ be a torsion section
and suppose $\tau_v^*\mathscr L$ and $\mathscr L$ are isomorphic
on the generic fiber.  Suppose the chosen generic-fiber lift
$A_v$ of translation satisfies $nv=0$ and
$A_v^n=\zeta$ as a line-bundle lift, for some $n\ge1$ and a
root of unity $\zeta$.
Then $A_v$ preserves the integral module of global sections and
its inverse does also.
\end{lemma}
\begin{proof}
Enlarge the finite base field if necessary so that the point,
the lift and the root of unity are defined.  Properness extends
the point to an integral section.  The generic line-bundle
isomorphism is a rational section of the line bundle
$\operatorname{Hom}(\tau_v^*\mathscr L,\mathscr L)$.
Its divisor has no horizontal component.  Since the special
fiber is geometrically integral and reduced, its only possible
vertical divisor is an integral multiple of that fiber.
Multiplying by a power of a uniformizer removes that divisor
and gives an integral isomorphism $B_v$.
Thus $A_v=cB_v$ for a nonzero scalar $c$.
The relation $A_v^n=\zeta$ includes $nv=0$, and
$B_v^n$ is an automorphism of the line bundle over the identity.
It is multiplication by a unit of $R$: all global functions on
the proper connected elliptic scheme are scalars.
Taking valuations in $c^nB_v^n=\zeta$ gives $v(c)=0$.
Hence $A_v$ and its inverse are integral, as required.
\end{proof}

The hypotheses here concern actual line-bundle lifts.  An arbitrary
finite-order matrix need not preserve an arbitrarily chosen
lattice; the vertical-divisor argument supplies the necessary
integral comparison.

\subsection{Norm preservation for finite combinations}

For a bounded power series $F$, write
$\|F\|=\sup_n|[T^n]F|_p$.

\begin{theorem}\label{ct-combination-norm}
Let $a\ge0$, let $s_1,\ldots,s_m,t$ be rational lattice points
whose orders modulo $\Lambda$ are prime to $p$, and let
$b_1,\ldots,b_m\in\mathbb C_p$ be arbitrary.  Then
\[
 \left\|\sum_{j=1}^m b_j
       \widehat{U_t(U_{s_j}\theta)}
                         (\omega\log(1+T))^a\right\|
 =
 \left\|\sum_{j=1}^m b_j
       \widehat{U_{s_j}\theta}
                         (\omega\log(1+T))^a\right\|.
 \tag{CT15}
\]
In particular every common coefficient bound is preserved.
\end{theorem}
\begin{proof}
The cases $a=0$ or $m=0$ are immediate.  Choose an even positive
integer $N$, prime to $p$, such that $Ns_j,Nt\in\Lambda$ for
all $j$.  On the complex elliptic curve $\mathbb C/\Lambda$,
the functions
\[
                           (U_{s_j}\theta)(Nz)^a
\]
are reduced sections of the same symmetric line bundle
\[
              \mathscr L_N=[N]^*\mathcal O([h])^{\otimes a}.
 \tag{CT16}
\]
To check the multiplier, translation by $s_j$ changes the
semicharacter on $N\Lambda$ by
$e^{2\pi iE(s_j,N\ell)}=1$, since $Ns_j\in\Lambda$.
The Hermitian form is $aN^2H$ for every summand.

Use the natural integral model of the divisor
$a[N]^{-1}[h]$.  Its support consists of nonzero torsion points
of order prime to $p$.  They do not reduce to the origin, because
the formal group has no nonzero prime-to-$p$ torsion.  Thus its
canonical section gives an integral trivialization near the
origin.  The corresponding reduced function is
$\theta(Nz)^a$.  By Proposition~\ref{ct-translates-integral}
and the integral substitution $T\mapsto(1+T)^N-1$, its
formal series is an integral unit.  Therefore the actual
reduced Taylor expansion of any section of (CT16) is its
expansion in the integral local trivialization multiplied by
this unit series.  Lemma~\ref{ct-expansion-saturated} identifies
its coefficient norm with the integral global-section norm,
also for arbitrary $\mathbb C_p$ linear combinations.

Put $v=t/N$.  Since $N^2v=Nt\in\Lambda$, translation by $v$
preserves the line bundle of degree $aN^2$.  In reduced analytic
coordinates its lift is
\[
       (A_vF)(z)=e^{-\pi aN^2 H(z,v)-\pi aN^2H(v,v)/2}
                                                        F(z+v).
 \tag{CT17}
\]
The algebraic Mumford translation has exactly this reduced
formula, as in the cited construction.  The lift is algebraic
after a finite extension.  By (CT3), a positive power of $A_v$
is the action of a lattice translation, which is a root of
unity because the line bundle is symmetric.  Lemma~\ref{ct-torsion-isometry}
therefore proves that it is an integral automorphism of the
global-section module.

Using $Nv=t$ in (CT17) gives the exact identity
\[
 A_v\bigl((U_{s_j}\theta)(Nz)^a\bigr)
                         =(U_t(U_{s_j}\theta))(Nz)^a.
 \tag{CT18}
\]
Hence the two sides of (CT15), after substituting $Nz$ for $z$,
have equal norm.  On the formal multiplicative group this
substitution is $T\mapsto(1+T)^N-1$; its inverse is
$T\mapsto(1+T)^{1/N}-1$, and both are integral with unit
linear coefficient.  They preserve the coefficient norm by
Lemma~\ref{ct-expansion-saturated}, or directly by applying
composition and its inverse to bounded series.  Removing this
substitution proves (CT15).
\end{proof}

This proves the entire finite-combination norm assertion for the
one-dimensional theta function, with coefficients allowed to be
arbitrary $p$-adic numbers and with one period fixed beforehand.
It does not rely on estimating the coefficients $b_j$ separately,
so cancellations in their sum are retained.

\subsection{Matching the chosen torsion character}

The multiplicative coordinate in (CT8) is unique only up to
$T\mapsto(1+T)^u-1$ with $u\in\mathbb Z_p^\times$.
The next argument makes this freedom explicit and matches the
character used for the selected CM torsion points.

\begin{proposition}\label{ct-character-normalization}
Let $p\ge5$ split in $k$, and let $\iota$ induce one of its
two primes.  For $n\ge0$, let $C_n$ consist of the classes
\[
 v=\frac{a+b\tau}{p^n}\pmod\Lambda,
 \qquad a,b\in\mathbb Z,\qquad
              \iota(a+b\bar\tau)\in p^n\mathbb Z_p.
 \tag{CT19}
\]
These are precisely the $p^n$-torsion points in the formal
subgroup.  The compatible maps
\[
  v\longmapsto\iota\!\left(e^{2\pi i b/p^n}\right)
      =\iota\!\left(e^{-2\pi iE(1,v)}\right)
 \tag{CT20}
\]
are isomorphisms $C_n\simeq\mu_{p^n}$.  One unit change of
the multiplicative coordinate makes its value $1+T(v)$ equal
to (CT20) simultaneously for every $n$.  The normalized period
$\Omega$ in (CT12) remains a unit.
\end{proposition}
\begin{proof}
Because $p$ splits and is unramified, both
$\iota(\tau)$ and $\iota(\bar\tau)$ lie in $\mathbb Z_p$,
and their difference is a unit.  For each residue of $b$ modulo
$p^n$, (CT19) determines one residue of $a$; hence $|C_n|=p^n$.
The formula for $E(1,v)$ follows from
$\operatorname{Im}\tau=\sqrt d/2$.  It shows that (CT20)
is a well-defined injective group map, hence an isomorphism.
These maps commute with multiplication by $p$.

To identify the formal subgroup, choose
$\beta\in\mathcal O_k$ with
$\iota(\beta)\equiv0\pmod{p^n}$ and
$\iota(\bar\beta)\equiv1\pmod{p^n}$.
In the multiplicative coordinate its action is
$1+T\mapsto(1+T)^{\iota(\beta)}$.
This identity follows first on formal logarithms, where the CM
endomorphism scales the invariant differential by
$\iota(\beta)$, and then on formal series in characteristic
zero.  Thus $\beta$ kills every formal $p^n$-torsion point.
For a complex torsion class represented by $(a+b\tau)/p^n$,
this forces $\beta(a+b\tau)\in p^n\mathcal O_k$, and its
conjugate component gives (CT19).  Conversely the formal
multiplicative group has $p^n$ such points, so cardinalities
give equality with $C_n$.

The initial multiplicative coordinate and (CT20) are now two
compatible isomorphisms from the same groups $C_n$ to
$\mu_{p^n}$.  Their ratio is an automorphism of the inverse
system, given by compatible units modulo $p^n$, hence by one
$u\in\mathbb Z_p^\times$.  Replacing $1+T$ by $(1+T)^u$
matches them.  Both this series and its inverse are integral,
and the period is divided by the unit $u$.  Its absolute value
is unchanged.
\end{proof}

In particular the source's stronger condition
$\iota(\bar\nu)\in p^{2n}\mathbb Z_p$ for a representative
$y=s\nu/p^n$, with $p\nmid s$, implies (CT19).  It also implies
$|y|^2=s^2N(\nu)/p^{2n}\in\mathbb Z$, a fact used in its
additional Weyl-translation phase.  The normalized chart above
matches the root of unity required there.  We now prove the
translation identity, including that phase.

\subsection{The exact translation by selected $p$-power torsion}

For $\zeta\in\mathbb C_p$ with $|\zeta-1|_p<1$ and a bounded
series $F(T)=\sum_m F_mT^m$, define
\[
 (S_\zeta F)(T)=F(\zeta(1+T)-1),\qquad
 [T^n]S_\zeta F
       =\zeta^n\sum_{m\ge n}\binom mn F_m(\zeta-1)^{m-n}.
 \tag{CT21}
\]
The displayed sums converge.  They give a bounded series with
$\|S_\zeta F\|\le\|F\|$, and the inverse operation is
$S_{\zeta^{-1}}$.  Thus $S_\zeta$ preserves the norm.
It preserves sums and products.  These statements follow directly
by evaluating on $|T|_p\le r<1$: every substituted argument has
absolute value at most $\max(r,|\zeta-1|_p)<1$, so the defining
series and all products converge uniformly there.  Alternatively
one can use (CT21) and the ultrametric convergence criterion.

We will also use the following compatibility.  If $R$ is a
rational function on $\mathscr E$, without poles in the formal
subgroup, and $v$ is a point of that subgroup with multiplicative
coordinate $\zeta$, then
\[
                 S_\zeta\widehat R=\widehat{R(z+v)}.
 \tag{CT22}
\]
The hats here use the same formal coordinate on both sides.
The boundedness required on the left follows from
Proposition~\ref{ot-elliptic-bounded}.  To justify the identity,
the integral formal-group parameter converges on the open unit
disc and its addition law there is the elliptic-curve addition
law.  Substituting its convergent coordinate functions into the
rational expression for $R$ gives the same analytic function
after translation by $v$.  Its denominators do not vanish in
this disc.  Expanding this identity at zero, using (CT21), proves
(CT22).  This argument uses translation of points on the elliptic
curve; it does not incorrectly replace a torsion point by its
formal logarithm, which is zero.

For the original even theta function put
\[
           (W_yF)(z)=e^{\pi i|y|^2}(U_yF)(z),\qquad
              \kappa(y)=e^{-2\pi iE(1,y)}.
 \tag{CT23}
\]

\begin{theorem}\label{ct-power-translation}
Let $s\ge1$ be prime to $p$, let $st\in\Lambda$, and let
\[
 y=\frac{s\nu}{p^n},\qquad \nu\in\Lambda,\qquad
                     \iota(\bar\nu)\in p^{2n}\mathbb Z_p.
 \tag{CT24}
\]
For the coordinate normalized in
Proposition~\ref{ct-character-normalization}, set
$\zeta=\iota(\kappa(y))$.  Then $\zeta^{p^n}=1$ and
\[
 \widehat{W_y(U_t\theta)}(\omega\log(1+T))
       =S_\zeta\!\left(
                \widehat{U_t\theta}(\omega\log(1+T))\right).
 \tag{CT25}
\]
In particular all series on the right and left are bounded,
with the same coefficient norm.
\end{theorem}
\begin{proof}
We first prove the assertion for $t=0$ and any $y=\nu/p^n$
with the stronger conjugate divisibility in (CT24).
For $n=0$, $y\in\Lambda$, and (CT25) is just (CT7), with
$\zeta=1$.  Suppose $n\ge1$.  Write $\mathfrak p$ for the
prime of $k$ induced by $\iota$.  Choose by the Chinese
remainder theorem
\[
 \epsilon\in2\mathcal O_k,\qquad
 \epsilon\equiv1\pmod{\mathfrak p^n},\qquad
 \epsilon\equiv0\pmod{\bar{\mathfrak p}^{\,n}}.
 \tag{CT26}
\]
It has nonzero unit image at $\mathfrak p$.
The function
\[
                     R(z)=\frac{\theta(\epsilon z)}
                                      {\theta(\bar\epsilon z)}
 \tag{CT27}
\]
is elliptic: the two numerator and denominator multipliers
coincide because $\epsilon$ and $\bar\epsilon$ have the same
norm, both in (CT6) and in the Hermitian exponential.
Its zeros and poles avoid the formal subgroup.  Indeed each
arises from a preimage of the nonzero point $h$ of order two,
whereas CM endomorphisms preserve the formal subgroup and that
subgroup has no nonzero prime-to-$p$ torsion.
Thus $R$ and its inverse have bounded expansions.

Put $f(T)=\widehat\theta(\omega\log(1+T))$.
For a CM element $\alpha$ denote its integral formal
endomorphism by
$[\alpha](T)=(1+T)^{\iota(\alpha)}-1$.
The formal expansion of (CT27) is
\[
                   \widehat R(T)=
                   \frac{f([\epsilon](T))}{f([\bar\epsilon](T))}.
 \tag{CT28}
\]
Both denominator and numerator have constant term one and are
integral.  Since $y$ belongs to the formal $p^n$-torsion,
\[
 [\epsilon](T\mathbin\oplus y)
     =[\epsilon](T)\mathbin\oplus y,\qquad
 [\bar\epsilon](T\mathbin\oplus y)=[\bar\epsilon](T),
\]
where $\oplus y$ means $T\mapsto\zeta(1+T)-1$.

The points $\bar\epsilon y$ and $(\epsilon-1)y$ belong to
$\Lambda$.  This can be checked at $\mathfrak p$ and
$\bar{\mathfrak p}$ from (CT24)--(CT26); all other primes
see no denominator.  Moreover $\bar\epsilon y\in2\Lambda$,
since $\epsilon\in2\Lambda$ and $p$ is odd.
The multiplier in (CT6) at $\bar\epsilon y$ therefore has
semicharacter one.  Cancelling the equal Hermitian exponential
factors gives the exact complex identity
\[
 R(z+y)=\frac{(U_{\epsilon y}\theta)(\epsilon z)}
                                      {\theta(\bar\epsilon z)}.
 \tag{CT29}
\]
Apply (CT22) to (CT28) and (CT29), then cancel the unit series
$f([\bar\epsilon](T))$.  Since $[\epsilon]$ is an integral
formal automorphism, substituting its inverse gives
\[
                      S_\zeta f=\widehat{U_{\epsilon y}\theta}.
 \tag{CT30}
\]
The formal variable on the right is still $\omega\log(1+T)$.

It remains to check the phase.  Write $\epsilon=A+B\tau$,
with $A,B\in\mathbb Z$, and put $l=(\epsilon-1)y\in\Lambda$.
The strong conjugate divisibility gives
$|y|^2=N(\nu)/p^{2n}\in\mathbb Z$.
From (CT3) and (CT6),
\[
 U_{\epsilon y}\theta
   =e^{\pi iE(y,l)-\pi i|l|^2}U_y\theta.
 \tag{CT31}
\]
Now
\[
 E(y,l)=-B|y|^2,\qquad
 |l|^2=(N(\epsilon)-2A-B+1)|y|^2.
\]
Therefore the exponent coefficient in (CT31) differs from
$|y|^2$ by
$(2A-2-N(\epsilon))|y|^2$, an even integer because
$\epsilon\in2\mathcal O_k$ makes $N(\epsilon)$ divisible
by four.  This proves
$U_{\epsilon y}\theta=W_y\theta$, and establishes (CT25)
for the base theta function.

Return to general $t$ and (CT24).  Put $y'=\nu/p^n$ and
$\zeta'=\iota(\kappa(y'))$, so $y=sy'$ and
$\zeta=(\zeta')^s$.  Let
\[
              q(z)=\frac{(U_t\theta)(z)}{\theta(z)},
                       \qquad Q(z)=q(sz).
 \tag{CT32}
\]
The function $Q$ is elliptic: its multiplier on translating by
$l\in\Lambda$ is $e^{2\pi iE(t,sl)}=1$ because
$st\in\Lambda$.  Its poles are among the preimages of $h$
under multiplication by $s$, so none lies in the formal subgroup.
It is regular at zero.  Proposition~\ref{ot-elliptic-bounded}
shows that $\widehat Q$ is bounded.  Since
$\widehat Q=\widehat q\circ[s]$ and $p\nmid s$, the inverse
of $[s]$ is integral; hence $\widehat q$ is bounded as well.

Apply (CT22) to $Q$ and $y'$.  The identity
$[s](T\mathbin\oplus y')=[s](T)\mathbin\oplus y$ and
the complex identity $Q(z+y')=q(sz+y)$ give, after cancelling
the invertible substitution $[s]$,
\[
                         S_\zeta\widehat q
                                     =\widehat{q(z+y)}.
 \tag{CT33}
\]
The right side is regular at zero: $y$ has $p$-power order and
therefore is not the theta zero of order two.
Finally the definition (CT23) gives
\[
                    W_y(U_t\theta)(z)=q(z+y)W_y\theta(z).
\]
Multiply (CT33) by the already proved base-theta identity and
use multiplicativity of $S_\zeta$.  This is exactly (CT25).
The norm assertion follows from (CT21); integrality of the
right side follows from Proposition~\ref{ct-translates-integral}.
\end{proof}

\subsection{What this classical route supplies}

For any given good ordinary CM model at the actual large split
primes, this section proves the precise matching of the even
reduced theta function, all prime-to-$p$ translated integrality,
one common discriminant-normalized unit period, preservation of
the coefficient norm of finite theta combinations, and the
normalization of its torsion character together with the exact
selected-$p$-power translation identity.  These are all three
one-dimensional theta laws used in the moment construction.
The exact higher-dimensional
lattice-theta construction and its finite translation/support
identities remain to be connected.  The arguments do not assert
an integral chart at a supersingular prime, do not treat $p=2$
or $3$ through the cited elliptic integrality statement, and do
not replace a reduced theta function by an arbitrary quadratic
exponential multiple.


